\part{Getting started}

% =====================================================================
\chapter{What CompLB3D does}
\label{ch:what}

CompLB3D simulates what happens in the pore space of a sediment, a soil or a
porous rock: water flows through it, dissolved chemicals move with the water,
microbes eat them and grow, minerals form and dissolve, and the pore space
itself changes shape as a result.

It does this at the \emph{pore scale}. The grid is fine enough that individual
grains and pore throats are resolved, so you are not fitting an effective
diffusivity or a tortuosity factor: the geometry is in the model, and transport
follows from it.

\section{The five things it couples}

\begin{center}
\begin{tikzpicture}[
  font=\footnotesize,
  box/.style={draw=cbRule, fill=cbSand, rounded corners=2pt, minimum width=3.1cm,
              minimum height=1.0cm, align=center, line width=0.7pt},
  arr/.style={-{Latex[length=2.4mm]}, cbGrey, line width=0.7pt}]

\node[box] (flow) at (0,0) {\textbf{Flow}\\ Navier--Stokes, D3Q19};
\node[box] (tr) at (4.2,0) {\textbf{Transport}\\ advection--diffusion, D3Q7};
\node[box] (rx) at (8.4,0) {\textbf{Reaction}\\ metabolism, chemistry};
\node[box] (eq) at (4.2,-1.9) {\textbf{Speciation}\\ equilibrium tableau};
\node[box] (geo) at (8.4,-1.9) {\textbf{Geometry}\\ precipitation, dissolution};

\draw[arr] (flow) -- node[above,font=\tiny]{velocity} (tr);
\draw[arr] (tr) -- node[above,font=\tiny]{concentrations} (rx);
\draw[arr] (rx) -- (eq);
\draw[arr] (eq) -- (tr);
\draw[arr] (rx) -- (geo);
\draw[arr] (geo.west) -- ++(-1.0,0) |- (flow.south);
\node[cbGrey,font=\tiny] at (2.0,-2.55) {pore space changed, flow re-solved};
\end{tikzpicture}
\end{center}

Every box except the first two is optional and switched on from the input file.
With all of them off, CompLB3D is a flow-and-transport solver and nothing more.
That is deliberate: you should be able to turn one thing on at a time and see
what it did.

\section{What makes it different from the 2D CompLaB}

CompLB3D began as a port of CompLaB to three dimensions and grew several
capabilities the 2D code does not have.

\begin{center}
\begin{tabular}{@{}p{5.2cm}p{4.6cm}p{4.6cm}@{}}
\toprule
& \textbf{CompLaB (2D)} & \textbf{CompLB3D}\\
\midrule
Dimensions & D2Q9 / D2Q5 & D3Q19 / D3Q7\\
Biomass solvers & CA, FD, LBM & CA, FD, LBM\\
Metabolism & kinetics, FBA, surrogate & the same, plus combined types\\
Aqueous speciation & --- & equilibrium tableau\\
Abiotic kinetics & --- & separate hook\\
Precipitation & --- & volume-of-pixel pore clogging\\
Dissolution & --- & pore reopening\\
Surrogate training & not distributed & MATLAB and Python paths\\
Worked examples & 6, one problem & 15, one per capability\\
\bottomrule
\end{tabular}
\end{center}

\section{What it is not}

Being honest about this saves time.

\begin{itemize}[leftmargin=1.2em]
\item \textbf{It is not a continuum reactive transport code.} There is no
Darcy law and no representative elementary volume. If your domain is metres
across, this is the wrong tool.
\item \textbf{The speciation is ideal.} No activity corrections, no temperature
dependence, no gas phase, no redox couples, no mineral saturation indices. The
equilibrium constants you supply are conditional constants at your own ionic
strength.
\item \textbf{There is no unsaturated flow.} One fluid phase only.
\item \textbf{Rate laws are compiled in.} Changing your chemistry means editing
a C\texttt{++} header and rebuilding. This surprises people who expect an input
file to be enough.
\end{itemize}

% =====================================================================
\chapter{Installing and building}
\label{ch:install}

\section{What you need}

\begin{center}
\begin{tabular}{@{}p{3.6cm}p{2.4cm}p{8.6cm}@{}}
\toprule
\textbf{Component} & \textbf{Required?} & \textbf{Notes}\\
\midrule
C\texttt{++} compiler & yes & GCC or Clang, C\texttt{++}11 or later\\
CMake & yes & 3.10 or later\\
Palabos & yes & \textbf{version 2.3.0}. Downloaded separately.\\
MPI & recommended & on by default; turn off with \code{-DENABLE\_MPI=OFF}\\
GLPK & optional & only for flux balance analysis through GLPK\\
Python 3 + cobrapy & optional & only for flux balance analysis through COBRApy\\
ParaView & recommended & to look at the output\\
\bottomrule
\end{tabular}
\end{center}

\warn{Palabos version matters.}{CompLB3D is developed against Palabos
\textbf{v2.3.0}. Later versions change data-processor interfaces, and the code
will not compile against them without edits. Get 2.3.0.}

\section{Step 1: get Palabos}

Palabos is a lattice-Boltzmann framework and it is a separate work under its
own licence. It is not distributed with CompLB3D, partly because it is around
100\,MB and partly because bundling someone else's AGPL library inside yours
muddies the licence question for everybody downstream.

Download \code{palabos-v2.3.0} from \url{https://palabos.unige.ch/} and unpack
it somewhere permanent.

\begin{lstlisting}[style=sh]
tar xzf palabos-v2.3.0.tar.gz -C /your/path/
\end{lstlisting}

You do not have to build Palabos separately --- CompLB3D compiles the parts it
needs. But building it once on its own is worth the ten minutes, because if
something is wrong with your compiler or MPI you will find out there, in a
project with good documentation, rather than here.

\begin{lstlisting}[style=sh]
cd /your/path/palabos-v2.3.0/build
cmake ..
make
\end{lstlisting}

Warnings during that build are normal and can be ignored as long as it finishes.

\section{Step 2: get CompLB3D}

\begin{lstlisting}[style=sh]
git clone https://github.com/shahram444/CompLB3D-lattice-Boltzmann-FBA-surrogate-COBRApy-GLPK.git
cd CompLB3D-lattice-Boltzmann-FBA-surrogate-COBRApy-GLPK
\end{lstlisting}

\section{Step 3: point CMake at Palabos}

Open \file{CMakeLists.txt} and edit the five lines that name the Palabos tree.
They are around lines 85 to 93.

\begin{lstlisting}[style=sh]
include_directories("/your/path/palabos-v2.3.0/src")
include_directories("/your/path/palabos-v2.3.0/externalLibraries")
include_directories("/your/path/palabos-v2.3.0/externalLibraries/Eigen3")
file(GLOB_RECURSE PALABOS_SRC "/your/path/palabos-v2.3.0/src/*.cpp")
file(GLOB_RECURSE EXT_SRC "/your/path/palabos-v2.3.0/externalLibraries/tinyxml/*.cpp")
\end{lstlisting}

This is the single most common reason a first build fails. If CMake reports
that it cannot find Palabos headers, this is why.

\section{Step 4: build}

\begin{lstlisting}[style=sh]
mkdir build && cd build
cmake ..
make -j
\end{lstlisting}

That gives you the plain build: flow, transport, hand-written kinetics,
speciation, precipitation and dissolution. It does \emph{not} include flux
balance analysis, because that needs an optional dependency.

\subsection{Optional: flux balance analysis through GLPK}

Install GLPK first --- your package manager will have it
(\code{apt install libglpk-dev}, \code{brew install glpk},
\code{module load glpk} on a cluster). Then:

\begin{lstlisting}[style=sh]
cmake -DENABLE_GLPK=ON ..
make -j
\end{lstlisting}

\subsection{Optional: flux balance analysis through COBRApy}

This embeds a Python interpreter, so you need the Python \emph{development}
headers, not just the interpreter.

\begin{lstlisting}[style=sh]
python3 -m pip install cobrapy
cmake -DENABLE_COBRAPY=ON ..
make -j
\end{lstlisting}

\note{Both at once.}{\code{cmake -DENABLE\_GLPK=ON -DENABLE\_COBRAPY=ON ..}
builds an executable that can do either, chosen at run time from the input
file. Worth doing once so you never think about it again.}

\subsection{The build options in full}

\begin{center}
\begin{tabular}{@{}p{4.6cm}p{1.6cm}p{8.6cm}@{}}
\toprule
\textbf{Option} & \textbf{Default} & \textbf{Effect}\\
\midrule
\code{ENABLE\_MPI} & ON & parallel build. Turn off for a single-core debug build.\\
\code{ENABLE\_GLPK} & OFF & links GLPK; enables \code{reaction\_type glpk}\\
\code{ENABLE\_COBRAPY} & OFF & embeds Python; enables \code{reaction\_type cobrapy}\\
\code{FBA\_BULK\_ONLY} & OFF & restricts the FBA processors to the bulk domain. An optimisation; leave it off until you have a reason.\\
\bottomrule
\end{tabular}
\end{center}

\section{Step 5: check the build}

The fastest confidence check is the example suite, which ships with the code.

\begin{lstlisting}[style=sh]
python3 examples/makeGeometry.py
python3 examples/makeExamples.py
python3 examples/validateExamples.py .
\end{lstlisting}

The last command checks the fifteen examples structurally --- tag names, vector
lengths, model file dimensions --- and compiles their kinetics headers. It
takes about a second and should end with:

\begin{lstlisting}[style=sh]
15 case(s) checked, 0 failed, 0 note(s).
\end{lstlisting}

Then run them:

\begin{lstlisting}[style=sh]
./examples/runAllExamples.sh .
\end{lstlisting}

\section{When the build fails}

\begin{center}
\begin{tabular}{@{}p{6.2cm}p{8.6cm}@{}}
\toprule
\textbf{Message} & \textbf{Cause and fix}\\
\midrule
\code{Cannot find palabos3D.h} & The paths in \file{CMakeLists.txt} are wrong. Step 3.\\
\code{ENABLE\_GLPK is ON but GLPK was not found} & Install \code{libglpk-dev}, or drop the flag.\\
\code{ENABLE\_COBRAPY is ON but the Python development headers were not found} & You have Python but not \code{python3-dev}. Install it.\\
undefined reference to \code{MPI\_*} & Building with MPI but linking without it. Load the same MPI module you configured with, and delete \file{build/} before reconfiguring.\\
errors inside Palabos headers & Almost always the wrong Palabos version. Check it is 2.3.0.\\
\code{defineKinetics.hh: C[94] out of range} & You are compiling a kinetics header written for different chemistry. Chapter~\ref{ch:ownkinetics}.\\
\bottomrule
\end{tabular}
\end{center}

\warn{Delete \file{build/} when you change configuration.}{CMake caches its
options. Switching \code{ENABLE\_GLPK} on and off without clearing the cache
produces builds that link the wrong things and fail in confusing ways.}

% =====================================================================
\chapter{Your first simulation}
\label{ch:first}

We will run example 01. It is the smallest thing that exercises both solvers:
water driven through a channel, and a dissolved tracer that rides along and
reacts with nothing.

\section{Run it}

\begin{lstlisting}[style=sh]
# from the top of the CompLB3D tree
cp examples/01_flow_only/defineKinetics.hh        .
cp examples/01_flow_only/defineAbioticKinetics.hh .
cd build && make -j && cd ..

mkdir -p run01 && cd run01
cp ../examples/01_flow_only/CompLaB.xml .
cp -r ../examples/01_flow_only/input .
mkdir -p output
../build/complab
\end{lstlisting}

\warn{Why the two \file{.hh} files get copied.}{CompLaB compiles its rate laws
in. \file{defineKinetics.hh} and \file{defineAbioticKinetics.hh} are
\code{\#include}d, not read at run time, so a case with different chemistry
needs a different binary. This catches everyone once. See
Chapter~\ref{ch:ownkinetics}.}

\note{\file{CompLaB.xml} is read by that exact name}{from the current working
directory. Not from the executable's directory, not from a path you pass on the
command line. Run the code from the folder that holds the input file.}

\section{What you should see}

The run prints a long banner at start-up. It is worth reading, because most
input mistakes are visible in it. The parts that matter:

\begin{lstlisting}[style=sh]
[OK] XML configuration loaded and validated

  domain          : 24 x 24 x 6   (3456 voxels)
  substrates      : 1
  microbes        : 0             (abiotic mode)
  ...
  [IMMOBILE] immobile (solid/precipitate) substrates: 0 of 1
\end{lstlisting}

then the flow solver converging, then the transport steps. At the end you have
a folder of \file{.vti} files in \file{output/}.

\section{Reading the run}

Three things are worth checking on every run you ever do.

\subsection{Did the flow converge?}

The Navier--Stokes solver iterates until the velocity field stops changing by
more than \tg{ns\_converge\_iT1}, or until \tg{ns\_max\_iT1} steps. If it hits
the iteration cap instead of the tolerance, your flow field is not converged
and everything downstream inherits that.

\subsection{Are there \code{[NEG!]} warnings?}

A negative concentration is not physical. It means a reaction removed more of
something than was there, which usually means the time step is too long for the
rate you specified. One or two at start-up while fields settle is tolerable;
a stream of them is a real problem. Chapter~\ref{ch:errors}.

\subsection{Does the answer depend on \textit{z}?}

Every shipped example is 24\,$\times$\,24\,$\times$\,6 with $z$ periodic and
the geometry extruded through the depth. That makes them quasi-2D on purpose:
every $z$-plane sees identical conditions, so the answer \emph{must not} depend
on $z$. If it does, something is wrong with the setup, not with the physics.
It is a free check and it costs nothing to look.

\section{Looking at the output}

CompLB3D writes VTK ImageData (\file{.vti}) files, one per field per save
interval. Open them in ParaView:

\begin{enumerate}[leftmargin=1.4em]
\item \textbf{File $\rightarrow$ Open}, select the numbered group
      (ParaView collapses \file{subsLattice0\_0000500.vti} and friends into one
      entry), and hit Apply.
\item Set the \textbf{Coloring} dropdown to the field you want.
\item Use \textbf{Slice} to cut through the domain; the raw view of a 3D block
      shows you only its outer surface.
\end{enumerate}

\note{The \file{.vti} files use unit spacing.}{ParaView therefore shows lattice
coordinates, not micrometres. Multiply by \tg{dx} if you need physical
distance. This is a known wart, not a setting.}

\section{What to try next}

Change one thing and rerun. Good first experiments:

\begin{center}
\begin{tabular}{@{}p{5.0cm}p{9.8cm}@{}}
\toprule
\textbf{Change} & \textbf{What should happen}\\
\midrule
\tg{Peclet} from 1 to 0 & flow switches off; the tracer profile becomes a straight line. This is example 02.\\
\tg{Peclet} from 1 to 10 & advection dominates; the tracer front sharpens and moves faster.\\
\tg{delta\_P} up & higher velocity, and the flow solver takes longer to converge.\\
\tg{nz} from 6 to 12 & the answer must not change. If it does, tell us.\\
\bottomrule
\end{tabular}
\end{center}

% =====================================================================
\chapter{The input file}
\label{ch:xml}

Everything about a run is in \file{CompLaB.xml}. This chapter walks through its
structure. Chapter~\ref{ch:reference} is the exhaustive tag list; this one is
the map.

\section{The six blocks}

\begin{center}
\begin{tikzpicture}[font=\footnotesize,
  b/.style={draw=cbRule,fill=cbSand,line width=0.7pt,minimum width=3.2cm,
            minimum height=0.85cm,align=center,rounded corners=2pt}]
\node[b] (p)  at (0,0)    {\texttt{<path>}};
\node[b] (s)  at (0,-1.1) {\texttt{<simulation\_mode>}};
\node[b] (l)  at (0,-2.2) {\texttt{<LB\_numerics>}};
\node[b] (c)  at (0,-3.3) {\texttt{<chemistry>}};
\node[b] (m)  at (0,-4.4) {\texttt{<microbiology>}};
\node[b] (o)  at (0,-5.5) {\texttt{<equilibrium>}, \texttt{<precipitation>},\\ \texttt{<dissolution>}, \texttt{<IO>}};

\node[anchor=west,cbGrey] at (2.1,0)    {where the input and output folders are};
\node[anchor=west,cbGrey] at (2.1,-1.1) {which physics is switched on};
\node[anchor=west,cbGrey] at (2.1,-2.2) {domain, geometry, time stepping};
\node[anchor=west,cbGrey] at (2.1,-3.3) {the solutes: one block each};
\node[anchor=west,cbGrey] at (2.1,-4.4) {the organisms: one block each};
\node[anchor=west,cbGrey] at (2.1,-5.5) {optional features and output};
\end{tikzpicture}
\end{center}

\section{\tg{simulation\_mode}: the master switches}

This block decides which physics runs at all. Everything in it is a
\code{true} or \code{false}.

\begin{lstlisting}[style=xml]
<simulation_mode>
    <biotic_mode>true</biotic_mode>
    <enable_kinetics>true</enable_kinetics>
    <enable_abiotic_kinetics>false</enable_abiotic_kinetics>
    <enable_fba_glpk>false</enable_fba_glpk>
    <enable_fba_cobrapy>false</enable_fba_cobrapy>
    <enable_surrogate>false</enable_surrogate>
</simulation_mode>
\end{lstlisting}

\tg{biotic\_mode} is the big one: \code{false} skips the whole
\tg{microbiology} block, and CompLB3D becomes a chemistry-and-transport code.

The three metabolic switches are \emph{global}. Each microbe then chooses its
own \tg{reaction\_type}, and the two have to agree --- a microbe asking for
\code{glpk} while \tg{enable\_fba\_glpk} is \code{false} stops the run with a
message. That double gate is deliberate: it means you cannot accidentally get
different biology from the one you asked for.

\section{\tg{LB\_numerics}: domain and time}

\begin{lstlisting}[style=xml]
<domain>
    <nx>24</nx> <ny>24</ny> <nz>6</nz>
    <dx>10</dx> <unit>um</unit>
    <characteristic_length>24</characteristic_length>
    <filename>geometry.dat</filename>
    <material_numbers>
        <pore>2</pore>
        <solid>0</solid>
        <bounce_back>1</bounce_back>
        <microbe0>3</microbe0>
    </material_numbers>
</domain>
\end{lstlisting}

\tg{dx} and \tg{unit} set the physical size of one voxel; here 10\,\textmu m.
\tg{characteristic\_length} is the length scale the Reynolds and P\'eclet
numbers are built on, in voxels.

\subsection{Material numbers}
\label{sec:materials}

Every voxel in the geometry file carries an integer. The
\tg{material\_numbers} block says what those integers mean.

\begin{center}
\begin{tikzpicture}[font=\footnotesize,scale=0.62]
\foreach \y in {0,...,7} \foreach \x in {0,...,11} {
  \pgfmathsetmacro{\v}{(\y==0 || \y==7) ? 1 : (((\x-3)*(\x-3)+(\y-4)*(\y-4) < 3) ? 0 : 2)}
  \pgfmathtruncatemacro{\vv}{\v}
  \ifnum\vv=1 \fill[cbGrey!55] (\x,\y) rectangle ++(1,1);\fi
  \ifnum\vv=0 \fill[cbRed!45] (\x,\y) rectangle ++(1,1);\fi
  \ifnum\vv=2 \fill[cbTeal!18] (\x,\y) rectangle ++(1,1);\fi
  \draw[white,line width=0.4pt] (\x,\y) rectangle ++(1,1);
}
\foreach \x in {7,8,9} \foreach \y in {1,2} {
  \fill[cbBlue!45] (\x,\y) rectangle ++(1,1);
  \draw[white,line width=0.4pt] (\x,\y) rectangle ++(1,1);}
\node[anchor=west,font=\scriptsize] at (12.6,7.0) {\tikz\fill[cbGrey!55](0,0)rectangle(0.3,0.3); \ \ \texttt{1} bounce\_back (walls)};
\node[anchor=west,font=\scriptsize] at (12.6,6.2) {\tikz\fill[cbRed!45](0,0)rectangle(0.3,0.3); \ \ \texttt{0} solid (grain)};
\node[anchor=west,font=\scriptsize] at (12.6,5.4) {\tikz\fill[cbTeal!18](0,0)rectangle(0.3,0.3); \ \ \texttt{2} pore (water)};
\node[anchor=west,font=\scriptsize] at (12.6,4.6) {\tikz\fill[cbBlue!45](0,0)rectangle(0.3,0.3); \ \ \texttt{3} microbe0 seed};
\node[anchor=west,font=\scriptsize,cbGrey,align=left] at (12.6,2.8)
  {one $z$-plane of a\\ 24 $\times$ 24 $\times$ 6 domain};
\node[cbGrey,font=\scriptsize] at (6,-0.7) {flow direction $\rightarrow$};
\end{tikzpicture}
\end{center}

The distinction between \code{solid} and \code{bounce\_back} matters: both stop
the flow, but \code{solid} voxels get no dynamics at all, while
\code{bounce\_back} voxels reflect the fluid. Use \code{bounce\_back} for the
domain walls and \code{solid} for grains.

A \tg{microbeN} entry does double duty: it says which voxels that organism is
seeded in, \emph{and} it marks the organism as a biofilm former rather than a
planktonic one. That has consequences described in Section~\ref{sec:biofilm}.

\subsection{The geometry file}
\label{sec:geomfile}

The geometry is a plain text file, one integer per line, $x$ fastest, then $y$,
then $z$:

\begin{lstlisting}[style=py]
for z in range(nz):
    for y in range(ny):
        for x in range(nx):
            write(mask[x][y][z])
\end{lstlisting}

So the file has exactly \code{nx*ny*nz} lines. There is no header. The examples
ship a generator, \file{examples/makeGeometry.py}, which is the easiest thing
to adapt for your own geometry.

\warn{Order, not shape.}{Nothing in the file records \code{nx}, \code{ny} or
\code{nz} --- those come from the XML. A file with the right number of values
but the wrong loop order produces a geometry that is silently transposed. If
your pore structure looks scrambled in ParaView, this is why.}

\section{\tg{chemistry}: the solutes}

One \tg{substrateN} block per solute, numbered from zero with no gaps.

\begin{lstlisting}[style=xml]
<chemistry>
    <number_of_substrates>2</number_of_substrates>

    <substrate0>
        <name_of_substrates>donor</name_of_substrates>
        <initial_concentration>0.1</initial_concentration>
        <substrate_diffusion_coefficients>
            <in_pore>5e-10</in_pore>
            <in_biofilm>5e-10</in_biofilm>
        </substrate_diffusion_coefficients>
        <left_boundary_type>Dirichlet</left_boundary_type>
        <left_boundary_condition>0.5</left_boundary_condition>
        <right_boundary_type>Neumann</right_boundary_type>
        <right_boundary_condition>0.</right_boundary_condition>
    </substrate0>
    ...
</chemistry>
\end{lstlisting}

\textbf{Dirichlet} holds the concentration at the boundary; \textbf{Neumann}
sets the flux, and \code{0} means no flux, i.e.\ a closed end.

Two optional flags are worth knowing early:

\begin{center}
\begin{tabular}{@{}p{4.0cm}p{10.8cm}@{}}
\toprule
\tg{immobile} & \code{true} means this substrate never advects or diffuses. It only receives its reaction source term. This is how minerals are represented: they must accumulate exactly where they form.\\
\tg{fix\_concentration} & a per-substrate list in \tg{chemistry}. Marks a solute as an inexhaustible reservoir: reactions read it but never draw it down.\\
\bottomrule
\end{tabular}
\end{center}

\section{\tg{microbiology}: the organisms}

One \tg{microbeN} block per organism. Two of its tags carry most of the
meaning, and they are independent of each other.

\begin{lstlisting}[style=xml]
<microbe0>
    <name_of_microbes>Bug</name_of_microbes>
    <solver_type>CA</solver_type>          <!-- how biomass MOVES -->
    <reaction_type>kinetics</reaction_type><!-- how it METABOLISES -->
    <initial_densities>1.0</initial_densities>
    <decay_coefficient>0.0</decay_coefficient>
    <viscosity_ratio_in_biofilm>0</viscosity_ratio_in_biofilm>
    <half_saturation_constants>0.05 0</half_saturation_constants>
    ...
</microbe0>
\end{lstlisting}

\subsection{Biofilm or planktonic}
\label{sec:biofilm}

An organism with a \tg{microbeN} entry under \tg{material\_numbers} is a
\emph{biofilm} former: it is seeded in those voxels and it can clog them. One
without is \emph{planktonic}: seeded uniformly through the pore space.

Three rules follow, and each stops a run when broken:

\begin{itemize}[leftmargin=1.2em]
\item \tg{viscosity\_ratio\_in\_biofilm} must be present for exactly the
biofilm organisms. The parser counts them and compares.
\item \tg{initial\_densities} must have one entry per material number the
organism is seeded in.
\item \code{solver\_type FD} is rejected for planktonic organisms.
\end{itemize}

\section{Where to look things up}

\begin{center}
\begin{tabular}{@{}p{6.4cm}p{8.4cm}@{}}
\toprule
\textbf{Question} & \textbf{Where}\\
\midrule
what does this tag do, exactly? & Chapter~\ref{ch:reference}, or the inline comments in \file{CompLaB\_reference\_template.xml}\\
what is a realistic value? & the example closest to your problem, Part~II\\
why did the run refuse to start? & Chapter~\ref{ch:errors}\\
\bottomrule
\end{tabular}
\end{center}

\note{\file{CompLaB\_reference\_template.xml}}{ships with the code and
documents every tag inline, next to a plausible value, with its units and known
limits. It is the fastest reference when you are actually editing a file.}
